% Abstract body only. No numerical results, as requested.
Bounded split optimization is an established response to myopic decision-tree induction, but its predictive gains must be assessed alongside its computational costs. We study single trees, pure bootstrap forests, and forests allocating different search horizons across members. Repeated paired partitions of a fixed classification benchmark compare depth-matched, unpruned learners at shallow, deeper, and unrestricted capacities, with and without feature subsampling. A separate training-only selection experiment applies shared structural grids to conventional and mixed-capable forest families. Sparse replacements improve specific shallow architectures, but the effects change with capacity and feature policy. The inclusive selected family has higher mean accuracy than selected random forests, alongside heterogeneous dataset effects, sensitivity to benchmark weighting, and substantially greater selection and refit costs. Its choices include pure forests, so this gain is not solely evidence for mixing. Mean-effect intervals and adjusted location tests are interpreted separately. The results support conditional benefits on a limited, previously explored sample, not a new induction algorithm, universal accuracy--cost superiority, or a validated per-dataset time-budget policy.
